%% Part VI: Doing it yourself. A working executor, and the mistakes to avoid.

\subsection{A hundred-line executor}
\label{sec:wakers-ch17}

\textit{Everything in this section fits into a program short enough to read in one sitting: a timer future that stores its waker, and an executor whose idea of ``wake'' is ``send me back to the queue.''}

Listing~\ref{lst:wakers-17.1} uses only the standard library. One simplification: its timer spawns a thread per timer to stand in for a driver, which keeps the example self-contained. Real runtimes never do this; they use the wheel of Section~\ref{sec:wakers-ch15}, and a waiting task never has a thread of its own (Lemma~\ref{lem:idle-cost-memory-only}).

\begin{lstlisting}[caption={A complete executor and timer.}, label=lst:wakers-17.1]
use std::future::Future;
use std::pin::Pin;
use std::sync::mpsc::{sync_channel, Receiver, SyncSender};
use std::sync::{Arc, Mutex};
use std::task::{Context, Poll, Wake, Waker};
use std::thread;
use std::time::Duration;

// ---------- A leaf future: a timer ----------
struct Shared {
    done: bool,
    waker: Option<Waker>,
}

pub struct Timer {
    shared: Arc<Mutex<Shared>>,
}

impl Timer {
    pub fn new(dur: Duration) -> Self {
        let shared = Arc::new(Mutex::new(Shared { done: false, waker: None }));
        let s = shared.clone();
        // Stand-in for a real reactor: a thread that "fires" later.
        thread::spawn(move || {
            thread::sleep(dur);
            let mut st = s.lock().unwrap();
            st.done = true;
            if let Some(w) = st.waker.take() {
                w.wake(); // "buzz the pager"
            }
        });
        Timer { shared }
    }
}

impl Future for Timer {
    type Output = ();
    fn poll(self: Pin<&mut Self>, cx: &mut Context<'_>) -> Poll<()> {
        let mut st = self.shared.lock().unwrap();
        if st.done {
            return Poll::Ready(());
        }
        // Not ready: remember WHO to wake. Always store the latest waker.
        st.waker = Some(cx.waker().clone());
        Poll::Pending
    }
}

// ---------- A tiny executor ----------
type BoxFut = Pin<Box<dyn Future<Output = ()> + Send>>;

struct Task {
    fut: Mutex<Option<BoxFut>>,
    queue: SyncSender<Arc<Task>>,
}

impl Wake for Task {
    // What "wake" means for THIS executor: put me back on the run queue.
    fn wake(self: Arc<Self>) {
        let q = self.queue.clone();
        q.send(self).expect("run queue full");
    }
}

fn spawn(f: impl Future<Output = ()> + Send + 'static, tx: &SyncSender<Arc<Task>>) {
    let task = Arc::new(Task {
        fut: Mutex::new(Some(Box::pin(f))),
        queue: tx.clone(),
    });
    tx.send(task).unwrap();
}

fn run(rx: Receiver<Arc<Task>>) {
    // Blocks (sleeps) when the queue is empty -- no busy spinning.
    while let Ok(task) = rx.recv() {
        let mut slot = task.fut.lock().unwrap();
        if let Some(mut fut) = slot.take() {
            let waker = Waker::from(task.clone());
            let mut cx = Context::from_waker(&waker);
            if fut.as_mut().poll(&mut cx).is_pending() {
                *slot = Some(fut); // park it until someone calls wake()
            }
        }
    }
}

fn main() {
    let (tx, rx) = sync_channel(100);
    spawn(async {
        println!("A: start");
        Timer::new(Duration::from_millis(300)).await;
        println!("A: woke up after 300ms");
    }, &tx);
    spawn(async {
        println!("B: start");
        Timer::new(Duration::from_millis(100)).await;
        println!("B: woke up after 100ms");
    }, &tx);
    drop(tx);
    run(rx);
    println!("all tasks done, executor exits");
}
\end{lstlisting}

Running it prints:

\begin{lstlisting}[language={}, numbers=none]
A: start
B: start
B: woke up after 100ms
A: woke up after 300ms
all tasks done, executor exits
\end{lstlisting}

Where each idea from this section appears in the listing:

\begin{itemize}
    \item \textbf{Leaf future (Definition~\ref{def:leaf-future}):} \code{Timer}, the only code that stores a waker.
    \item \textbf{Clone to keep (Lemma~\ref{lem:clone-to-keep}):} \code{st.waker = Some(cx.waker().clone())}.
    \item \textbf{Atomic check-and-store (Lemma~\ref{lem:atomic-check-and-store}):} the \code{done} check and the waker store happen under the same lock that the timer thread takes before calling the waker.
    \item \textbf{Waker slot (Definition~\ref{def:waker-slot}):} the \code{waker} field of \code{Shared}.
    \item \textbf{Raw waker (Definition~\ref{def:raw-waker}):} built safely by implementing \code{Wake} for \code{Task} and calling \code{Waker::from(task.clone())}.
    \item \textbf{Waker-driven executor (Definition~\ref{def:waker-driven-executor}):} \code{run} polls only tasks received from the queue, and \code{wake} is the only way back onto it.
    \item \textbf{Remote wake:} the timer thread calls the waker; the executor thread is blocked in \code{recv()}, which costs no CPU (Axiom~\ref{ax:free-blocking}).
    \item \textbf{Independence:} the timer and the executor know nothing about each other. The only thing they share is the \code{Waker}.
\end{itemize}

When the last task finishes, every sender to the queue has been dropped, \code{recv()} fails, and the loop exits.

\subsection{Common mistakes}
\label{sec:wakers-ch18}

\begin{table}[htbp]
    \centering\small
    \begin{tabularx}{\linewidth}{>{\raggedright\arraybackslash}p{4.2cm} Y >{\raggedright\arraybackslash}p{1.8cm}}
        \toprule
        \textbf{Belief or habit} & \textbf{What is actually true} & \textbf{See} \\
        \midrule
        ``\texttt{wake()} runs the task.'' & It only schedules it. The poll happens later, on whichever thread the executor chooses. \texttt{wake()} must be cheap and must not block. & Def.~\ref{def:waker} \\
        ``Each waker is a background thread.'' & It is two pointers in a waker slot. Idle cost is memory only. & Lemma~\ref{lem:idle-cost-memory-only} \\
        ``The kernel calls my waker.'' & The kernel returns from a system call; Rust code calls the waker. & Lemma~\ref{lem:wakers-called-from-user-space} \\
        Storing only the first waker & Store the latest on every \texttt{Pending}; use \texttt{will\_wake} to avoid needless clones. & Axiom~\ref{ax:wake-obligation} \\
        Checking, then storing, separately & A lost wake can fall between the two. & Lemma~\ref{lem:atomic-check-and-store} \\
        Calling \texttt{wake\_by\_ref()} on every poll & Occasionally a deliberate yield; every time, a busy loop. & Table~\ref{tab:wakers-9.1} \\
        Blocking inside async code & \texttt{std::fs::read} or \texttt{thread::sleep} in a task blocks the whole worker. Use \texttt{tokio::fs}, \texttt{tokio::time::sleep} or \texttt{spawn\_blocking}. & Section~\ref{sec:wakers-ch16} \\
        ``It works inside \texttt{join!}.'' & Possibly an accidental rescue. & Lemma~\ref{lem:shared-fate} \\
        Spawning onto \texttt{current\_thread} from elsewhere & Nothing runs unless some thread is in \texttt{block\_on}. & Cor.~\ref{cor:no-block-on-no-progress} \\
        Testing a future by hand & \texttt{Waker::noop()}, stable since Rust 1.85, provides a do-nothing waker for polling once in a unit test. & --- \\
        \bottomrule
    \end{tabularx}
    \caption{Misconceptions and their corrections.}
    \label{tab:wakers-18.1}
\end{table}

\begin{Exercise}[title={Building an executor}]
    \begin{enumerate}
        \item In Listing~\ref{lst:wakers-17.1}, suppose \code{Timer::poll} stored the waker only on its first poll (\code{if st.waker.is\_none()}). Does this program still work? When would that change break a future?
        \item Why does \code{run} return once both tasks have finished? What happens if \code{main} keeps a clone of \code{tx} alive?
        \item Write a \code{yield\_now()} future for this executor: it returns \code{Pending} once, asking to be polled again straight away, and \code{Ready} the second time.
        \item Write a unit test that polls a fresh \code{Timer} once, without any executor, and asserts that it is pending.
    \end{enumerate}
\end{Exercise}
\begin{Answer}
    \begin{enumerate}
        \item This program still works: every poll of a task is given a waker for the same \code{Arc<Task>}, so the first waker stays valid. It breaks the Wake Obligation, though, which refers to the waker from the most recent poll (Axiom~\ref{ax:wake-obligation}). It fails when the waker changes between polls, for example inside \code{FuturesUnordered} or \code{select!} implementations that wrap the waker, when the future is moved into another task, or on runtimes whose wakers identify the worker. Store the latest waker, using \code{will\_wake} to skip needless clones.
        \item Each \code{Task} holds a \code{SyncSender}, and \code{main} drops its own \code{tx}. When a task finishes, its future is dropped and, once the last \code{Arc<Task>} (including any waker) is gone, so is its sender. With every sender dropped, \code{recv()} returns an error and the loop ends. If \code{main} kept a \code{tx} alive, \code{recv()} would block forever after the last task: the program would hang, using no CPU.
        \item The future wakes itself before returning \code{Pending} (the self-wake of Section~\ref{sec:wakers-ch7}):
\begin{lstlisting}
struct YieldNow { yielded: bool }

impl Future for YieldNow {
    type Output = ();
    fn poll(mut self: Pin<&mut Self>, cx: &mut Context<'_>) -> Poll<()> {
        if self.yielded {
            return Poll::Ready(());
        }
        self.yielded = true;
        cx.waker().wake_by_ref(); // back onto the run queue right away
        Poll::Pending
    }
}

fn yield_now() -> YieldNow { YieldNow { yielded: false } }
\end{lstlisting}
        The wake arrives during the poll, while \code{run} still holds the task's \code{fut} lock. That is fine: \code{wake} only sends the \code{Arc<Task>} on the channel, and \code{run} receives it after the current iteration has put the future back and released the lock.
        \item \code{Waker::noop()}, stable since Rust 1.85, gives a waker that does nothing; \code{Timer} is \code{Unpin}, so \code{Pin::new} works:
\begin{lstlisting}
#[test]
fn fresh_timer_is_pending() {
    let mut timer = Timer::new(Duration::from_secs(60));
    let mut cx = Context::from_waker(Waker::noop());
    assert!(Pin::new(&mut timer).poll(&mut cx).is_pending());
}
\end{lstlisting}
    \end{enumerate}
\end{Answer}

\subsection{The whole story in one page}
\label{sec:wakers-afterword}

A worker thread takes a task from its run queue and polls it, handing over a \code{Context} that holds the task's \code{Waker}. Through each \code{.await}, the context travels down the task's tree of futures to the leaf that is actually waiting. If the leaf can't finish, it leaves a clone of the waker in a waker slot---a socket's reader slot, a slot in the timer wheel, a blocking job's handle---and returns \code{Pending}, which bubbles back up to the worker.

When the worker runs out of tasks, it goes on driver duty. It asks the timer wheel for the next deadline and sleeps in mio's \code{Poll::poll}, which on Linux is \code{epoll\_wait}, with that deadline as the timeout. Three things can end the sleep:

\begin{itemize}
    \item the kernel reporting that a socket is ready;
    \item the timeout expiring;
    \item another thread performing a thread-level wake through \code{mio::Waker}.
\end{itemize}

In each case, ordinary Rust code then calls the right task's \code{wake()}. The task goes back onto a run queue, is polled again, and this time the leaf returns \code{Ready}, and the \code{async fn} continues where it paused.

Two facts summarize the rest. If nobody ever calls that waker, the task is stranded and never polled again (Lemma~\ref{lem:stranding}). And because a waiting task's idle cost is memory only (Lemma~\ref{lem:idle-cost-memory-only}), a single thread can hold hundreds of thousands of them.
